\chapter{Capstone 应急演练与替班教师接管手册}

\section{案例导读：替班接管测试的是课程包的独立性}

如果一门课只有原教师本人能解释清楚，那么它还不是稳定课程包。临时缺席、助教变动或项目中断都会暴露一个问题：别人能不能在有限时间内接手，并保护学生进度不被打断。

本章把 substitute handoff 看成一次应急演练。它要求把关键入口、当前状态、未决风险、评分边界和沟通模板提前准备好，让课程包在人员变化时仍然保持可运行。

\section{案例目标}

第 54 章把课程运行中的故障模式、暂停条件和升级边界整理成 escalation playbook。本章继续处理一个更现实、也更容易被忽视的场景：\textbf{当主讲教师临时缺席、机房环境出故障、助教临时短缺，或者课程必须由替班教师最小化接手时，团队怎样用一份 substitute handoff package 维持课程连续运行。}

本章的目标有四点：

\begin{enumerate}
    \item 理解 contingency drill 不是悲观预设，而是让课程在人员变化下仍能继续运行；
    \item 学会检查 trigger、minimal run package、substitute owner、context brief 和 escalation 状态；
    \item 学会区分“替班可接手”“补最小运行包”“先指定替班 owner”“刷新上下文 brief”和“先升级再替班”五类出口；
    \item 学会把课程从“依赖某个人记得细节”推进到“任何接班人都能在压力下安全接手”的状态。
\end{enumerate}

\section{最脆弱的时刻往往是临时缺席}

课程正常运行时，很多隐含知识都藏在主讲教师和熟悉课程的助教脑子里：哪一段 notebook 最容易出错，哪一项评分标准最容易被误解，哪条数据权限说明必须当面讲清，哪一周最容易被进度拖慢。

一旦主讲教师请假、助教临时短缺、机房环境失效，或者突然需要别的教师代上一节，课程最危险的不是“材料没有”，而是接班人拿到的东西虽然很多，却不知道今天最少要做什么、哪些边界不能碰、哪些情况必须立刻升级。

所以，本章的核心立场是：\textbf{替班 handoff 的目标不是完整复刻主讲教师，而是提供一个在压力下仍可安全执行的 minimal run package。}

\section{一个极小的 substitute-handoff 数据集}

配套 notebook 使用 \nolinkurl{capstone_contingency_substitute_handoff_demo.csv}。这是一组完全本地、教学用途的\textbf{应急演练与替班接管卡片}。每一行代表一个可能触发替班或应急接手的场景，包含：

\begin{itemize}
    \item \texttt{handoff\_id}；
    \item \texttt{contingency\_area}；
    \item \texttt{trigger\_status}；
    \item \texttt{minimal\_run\_package\_status}；
    \item \texttt{substitute\_owner\_status}；
    \item \texttt{context\_brief\_status}；
    \item \texttt{escalation\_status}；
    \item \texttt{readiness\_score}；
    \item \texttt{reference\_route}。
\end{itemize}

这里的 route 分成五类：

\begin{itemize}
    \item \texttt{substitute\_ready}：替班教师或替班助教可以按 minimal package 接手；
    \item \texttt{complete\_minimal\_run\_package}：最小运行包还不完整；
    \item \texttt{assign\_substitute\_owner}：还没有明确的接班 owner；
    \item \texttt{refresh\_context\_brief}：上下文说明还是旧版本或不够清楚；
    \item \texttt{escalate\_before\_substitution}：触及政策、隐私或关键边界，必须先升级。
\end{itemize}

\section{先看一个危险 baseline：只按 readiness score 排序}

当课程团队准备替班方案时，一个常见 baseline 是只按 readiness score 排序：看起来最成熟的场景优先放进手册。

\begin{examplebox}[只按 readiness 选择替班方案]
\begin{lstlisting}[style=pythonstyle]
top4 = sorted(
    rows,
    key=lambda row: row["readiness_score"],
    reverse=True,
)[:4]
\end{lstlisting}
\end{examplebox}

这个 baseline 很容易把“看起来差不多”的场景和真正能安全接手的场景混在一起。尤其是当 minimal run package 还不完整、context brief 过期，或者某个场景本身带有政策边界时，高分反而会制造过度自信。

在本章教学数据上，只按 readiness score 取前四项，真正替班可接手的精度只有 0.25，未就绪场景混入率是 0.75。表~\ref{tab:ch55_readiness_vs_substitute} 展示了结构化 substitute workflow 如何先看 escalation，再看 owner，再检查 minimal run package，最后确认 context brief：

\begin{table}[htbp]
\centering
\small
\setlength{\tabcolsep}{4pt}
\begin{tabular}{@{}p{0.28\textwidth}>{\centering\arraybackslash}p{0.18\textwidth}>{\centering\arraybackslash}p{0.18\textwidth}p{0.28\textwidth}@{}}
\toprule
方法 & 替班可接手精度 & 未就绪混入率 & 教学解读 \\
\midrule
只按 readiness 取前四项 & 0.25 & 0.75 & 分数高不代表接班人真能在压力下安全接手 \\
结构化 substitute workflow & 1.00 & 0.00 & 先判边界，再定人，再补最小运行包和上下文 \\
\bottomrule
\end{tabular}
\caption{本章 notebook 中两个最小替班流程的对照。替班接管不是 readiness 排名，而是安全可执行状态表。}
\label{tab:ch55_readiness_vs_substitute}
\end{table}

\section{一个可执行的 substitute workflow}

本章 notebook 的 routing 逻辑可以写成：

\begin{examplebox}[一个最小替班接管 routing 逻辑]
\begin{lstlisting}[style=pythonstyle]
if escalation_status in {"policy", "privacy", "grading"}:
    escalate_before_substitution
elif substitute_owner_status != "assigned":
    assign_substitute_owner
elif minimal_run_package_status != "ready":
    complete_minimal_run_package
elif context_brief_status != "ready":
    refresh_context_brief
else:
    substitute_ready
\end{lstlisting}
\end{examplebox}

这个顺序体现了一个应急原则：\textbf{边界先于替班，人先于材料，最小运行包先于完整课程体验。}

\section{替班手册应该包含什么}

一份可执行的 substitute handoff package 至少应该包含：

\begin{enumerate}
    \item 今天最少要完成什么：本节课最小教学目标和不可跳过的 gate；
    \item 材料入口：学生 handout、notebook、slides、rubric、签到或提交入口；
    \item 危险边界：哪些问题可以现场处理，哪些必须升级；
    \item 课堂节奏：90 分钟或 120 分钟内的最小流程；
    \item 接班联系人：谁能回答内容问题、谁能处理政策问题、谁能处理技术问题；
    \item 事后交回记录：替班时发生了什么、有哪些临时变更、哪些问题留到下次修。
\end{enumerate}

它不应该是一本厚手册，而应该是接班人临时接到任务时也能快速执行的一页或两页 briefing。

\section{配套 notebook 做了什么}

本章 notebook 采用的是一个 teaching-first 的最小设计：

\begin{itemize}
    \item 先读取 substitute-handoff table，并统计 trigger、package、owner、brief 和 reference route；
    \item 用 readiness score 做一个 naive substitute baseline；
    \item 再实现一个带 escalation、owner、minimal package 和 context brief 的 substitute workflow；
    \item 输出 substitute-ready、complete-package、assign-owner、refresh-brief 和 escalate-first 五个队列；
    \item 为每个非 ready 场景生成替班前 checklist；
    \item 最后生成 substitute handoff outline 和 contingency assistant prompt。
\end{itemize}

\section{AI 助手如何帮助你}

在这个案例里，AI 助手最适合帮助你：

\begin{itemize}
    \item 把完整课程包压缩成替班教师能读懂的 minimal run package；
    \item 从上学期记录中提炼“今天最少要做什么”的短 briefing；
    \item 标出哪些情况仍然必须升级给主讲教师或课程负责人；
    \item 把复杂上下文改写成替班教师上课前 5 分钟能读完的 handoff note；
    \item 帮你生成 contingency drill 演练模板。
\end{itemize}

但 AI 助手不能替团队决定是否允许替班触碰政策或评分边界，也不能替代课程负责人的最终授权。它能让接班更清楚，却不能替你承担课堂责任。

\section{练习}

\begin{exercisebox}[基础练习]
把一个 \texttt{substitute\_ready} 场景的 \texttt{context\_brief\_status} 改成 \texttt{stale}。重新运行 notebook，观察它为什么要先离开 ready 队列。
\end{exercisebox}

\begin{exercisebox}[进阶练习]
新增一个 readiness 很高、但 \texttt{minimal\_run\_package\_status = partial} 的替班场景。预测它会落到哪个 route，并说明为什么“看起来差不多”不等于“替班可执行”。
\end{exercisebox}

\begin{exercisebox}[开放练习]
为你自己的课程项目写一页 substitute handoff brief。至少包含：最小教学目标、材料入口、危险边界、接班联系人、课堂节奏和事后交回记录格式。
\end{exercisebox}

\section{本章小结}

课程团队真正成熟的标志之一，是关键成员缺席时课程仍能被安全接住。本章把 trigger、minimal run package、substitute owner、context brief、escalation status 和 readiness 放进同一张 handoff card。

当课程团队能把场景分成替班可接手、补最小运行包、先指定替班 owner、刷新上下文说明和先升级再替班，Capstone 就从“依赖某一个人撑住”走向“团队可替补、课程可持续”。
